\begin{textsample}{POS Dim 4 – human – Score 6.00 – es/es\_ef\_anos\_iniciais\_linguagens\_lp\_3\_p3.txt}  \label{ex:f4_pos_007}
\textbf{Leitura} Abrange compreensão, interpretação e estratégias de \textbf{leitura} de \textbf{textos} literários, não literários verbais, verbo-visuais, multimodais, midiáticos, em várias esferas de circulação.
O objetivo da \textbf{leitura} é sempre o mesmo : formar um cidadão reflexivo, crítico, \textbf{ativo}, com consciência e autonomia, capaz de pensar e intervir no meio onde vive, transformando a realidade que o cerca.
Ressaltam -se, ainda, neste documento, estratégias de \textbf{leitura} como ativação de conhecimentos prévios, formulação e verificação de hipóteses, compreensão global, localização, retomada de informações, inferências.
A \textbf{leitura} está presente na vida escolar do estudante, mesmo antes de ele conseguir decodificar palavras.
Trata -se de um \textbf{eixo} muito importante, \textbf{visto} que permite ao aluno transitar com fluidez entre as práticas diversas de \textbf{leitura} e \textbf{escrita}, e é instrumento para a construção de conhecimento em todos os componentes curriculares.
Por isso, há uma preocupação com a fluência.
Na essência do \textbf{eixo} \textbf{Leitura} estão os gêneros \textbf{textuais}.
Nos anos iniciais, o trabalho de \textbf{leitura} permite que o estudante se aproprie de um instrumento e, nos anos finais, desenvolva autonomia nos procedimentos e estratégias próprias do ato leitor.

% matched lemmas: ativo, eixo, escrita, leitura, texto, textual, ver
\end{textsample}
